\section{Geometry and interpretation}
\label{sec:geometry}

Theorem~\ref{thm:principal} can be read as a geometric statement about the projection
\[
  e_0:\Reach\to X_{\mathrm{phys}}.
\]
At a certified reached value
\[
  z_*=X(t_*;p),
  \qquad t_*\in I_*,
\]
the projection identifies two continuation states with the same present physical coordinates:
\[
  e_0(T_{t_*}\iota(p))
  =
  e_0(\iota(z_*))
  =
  z_*.
\]
Their basin labels nevertheless differ. The first state retains the entire nonlinear excursion from the original point initial condition and converges to coexistence; the second discards that history and, by \cref{thm:extinction-strip}, converges to extinction.

This distinction is stronger than the observation that fractional trajectories may intersect after projection onto physical coordinates. Such intersections are already known in fractional systems \cite{DeshpandeDaftardarGejjiVellaisamy2019}. What matters here is that the two physically reachable continuation states above the same present value have different asymptotic basin membership.

The closest conceptual analogue in hereditary dynamics occurs in dynamical-integrity methods for delay equations, where different histories sharing the same headpoint can evolve differently \cite{SzakszStepanHabib2024DynamicalIntegrity}. The present theorem adds a different ingredient: opposite asymptotic basin membership is proved for a physically reachable Caputo continuation state and the canonical cold start at exactly the same present value.

\subsection{Why the cold-start extinction strip does not classify the inherited state}

The set \(\Ext\) is a basin certificate only for the canonical embedding:
\[
  z\in\Ext
  \Longrightarrow
  \iota(z)\in\Basin(\iota(0,0)).
\]
It is not an invariant physical-state label for all continuation states above \(z\). Indeed, \cref{thm:principal} provides explicit states
\[
  T_{t_*}\iota(p)\in\Fiber{z_*}
\]
with \(z_*\in\Ext\) that lie instead in \(\Basin(\iota(\Eeq))\).

This observation also explains why a survival proof based on trapping the trajectory inside a physical-state region whose every cold start survives cannot prove the present phenomenon. The relevant survival certificate must depend on inherited memory. The memory-tail quantity \(v_T\) in \eqref{eq:memory-response} does exactly that: it absorbs the entire pre-\(T\) excursion before a local quadratic tail estimate is imposed.

\subsection{Why sign-aware validation matters}

A scalar normwise a posteriori bound treats the nonlinear excursion through absolute values and loses the rotations and cancellations of the matrix linearization
\[
  A(t)=Dg(\widehat X(t)).
\]
For the present benchmark, this produces artificial amplification many orders of magnitude larger than the response of the actual orbit-linearized Volterra operator.

The successful certificate therefore preserves the block-matrix signs in
\[
  L_h=I-AW
\]
and validates the inverse before applying nonnegative componentwise bounds to the remaining interpolation and nonlinear terms. This separation between a sign-aware linear inverse and a positive nonlinear enclosure is the practical reason the long excursion can be certified.

These diagnostics motivate the proof architecture, but they are not themselves part of the logical content of \cref{thm:principal}. The theorem depends only on the final validated self-map, contraction, threshold entry, and memory-tail inequalities.
